\section{Level symmetries and the original generic descent}
\label{sec:pel-level-symmetry}

The local integral model can be chosen to preserve the original finite
level symmetries.  Its generic compactification also descends to the
original smaller field when the full auxiliary level is one.  These
are actual clauses of the geometric construction; they do not assert
the whole simultaneous level-system contract.

Put $L=\operatorname{lcm}(M,A)$ and let $U(L)$ be full principal
level in the integral PEL similitude group $K_0=G(\widehat{\mathbb Z})$.
Choose the projective cone decomposition in
Theorem~\ref{thm:original-integral-pel-model} to be $K_0$-invariant.
This is possible because $U(L)$ is normal in $K_0$, by
\cite[Proposition 2.8(2)]{level-symmetry-lan-cones}.
Use the singleton original lattice throughout.

\begin{proposition}[Extending the prescribed level actions]
\label{prop:pel-original-level-actions}
The resulting local model satisfies the source's
\texttt{LevelActs} condition.  The extension preserves the actual
semiabelian family, coefficient action and rigidified primitive line.
\end{proposition}
\begin{proof}
Start with an automorphism $r$ of the original generic moduli space,
an isomorphism of its original semiabelian line data over $r$, and
the two fixed coefficient-linear label maps in the definition of
\texttt{LevelActs}.  Complete the half-labels using the prescribed
root, as in Theorem~\ref{thm:original-half-level-completion}.
The two label maps then determine one automorphism of the full
$L$-level module: on the missing half it is forced by the perfect
cross-pairing, and the overlap with $A$-division is unchanged.
This automorphism preserves the $L$-pairing root.

It is liftable to $K_0$.  At a geometric moduli point compare
coherent Tate markings before and after the given family
isomorphism.  Their ratio is an integral PEL similitude, and its
reduction is the just-determined finite-level automorphism.  This
uses the original full-Tate condition; arbitrary surjectivity of
an integral group onto its reduction is not being assumed.
The ratio may be chosen at any nonempty component, since the
finite-level action has already been fixed by the two given maps.
If the moduli space is empty, its dense integral model is empty
and the assertion is immediate.

Thus $r$ is the action of this element of $K_0/U(L)$.  Indeed the
two actions give the same underlying object and all labels, and
the original fine moduli property and automorphism rigidity give
equality.  The same uniqueness identifies the specified generic
family isomorphism.

The equivariant toroidal construction extends this action when
the cones are invariant
\cite[Section 7, especially Proposition 7.3]{level-symmetry-lan-cones}.
For an element of $K_0$ the lattice is unchanged: this is a change
of labels, with no isogeny of the abelian family.  The fixed-root
component is preserved because its multiplier is one modulo $L$.
The action extends to its reduced closure and then to its finite
normalization by functoriality of normalization.  These extensions
are inverse to those of the inverse group element.

The extended tautological family has the original coefficient
action.  Full faithfulness for cubical lines over the normal base
extends the original line comparison uniquely, including its
rigidification, adjunction and primitive identities.  The resulting
automorphism and family map satisfy every equality in
\texttt{LevelActs}, since those equalities hold on the dense
generic open.  No extra scalar is chosen.
\end{proof}

\begin{proposition}[Descent of the generic compactification]
\label{prop:pel-original-generic-descent}
If $A=1$, the same choices can be made so that the local model
satisfies the original \texttt{Descends} condition.  In particular
its generic compactification and semiabelian family descend to
the source field $K_A$, while their pullbacks give exactly the
original objects over the prescribed coefficient field.
\end{proposition}
\begin{proof}
Let $H$ be the original half-level stabilizer.  The completed-level
calculation gives
\[
 H_0=H\cap\{\nu\equiv1\pmod M\}=U(M).
\]
The multiplier induces an isomorphism
$H/H_0\simeq(\mathbb Z/M\mathbb Z)^\times$.
Injectivity was proved in Corollary~\ref{cor:original-neat-level}.
For surjectivity, choose a profinite lift of a residue unit.  At
the split places above $M$ use the coefficient unit $(1,r)$,
which fixes the half-marking and has norm $r$; at the other places
use the unramified norm on units.  This constructs an element of
$H$ with the required multiplier, over the same original lattice.

Consider the full-$M$ toroidal compactification over $K_A$, keeping
all primitive pairing roots, with the invariant cones chosen above.
Its root map to the finite \(K_A\)-scheme of primitive $M$-th
roots extends across the boundary: roots are integral, and the
compactification is normal.  The action of $H/H_0$ on these roots
is simply transitive on every geometric fibre.  It is therefore
free on the compactification, including its boundary.

The finite quotient is a scheme and the quotient morphism is
finite \'etale.  One can form it using an invariant ample bundle:
tensor the translates of an ample bundle on the projective
compactification.  The same equivariant construction carries the
tautological semiabelian family.  Its symmetric cubical line is
relatively ample: apply the generic ampleness criterion
\cite[Lemma 4.2.1.6]{level-symmetry-lan} after completion at each
normal local base ring, and descend relative ampleness by faithful
flatness.  Hence the family is
quasi-projective locally on the base and its equivariant descent
is effective in schemes.  All group, coefficient and cubical
structures descend by finite \'etale descent.

On the dense interior this quotient is the original half-level
moduli object over $K_A$.  Over the prescribed coefficient field,
the already chosen root $\zeta_M$ selects one member of each
$H/H_0$-orbit.  Consequently its fixed-root compactification is
canonically the base change of the quotient.  The same statement
holds for the entire semiabelian family and its comparison on the
interior.  This gives the objects, cartesian squares and original
generic family equalities required by \texttt{Descends}.
The coefficient field need not be Galois over $K_A$.
\end{proof}

\begin{proposition}[The elementary extension clauses]
\label{prop:pel-prime-to-p-and-dvr-extension}
On the abelian locus of the actual local model, every prescribed
generic torsion label of order prime to $p$ extends uniquely.
The model satisfies the original \texttt{GoodRedExt} clause over
every valuation ring with the specified fraction field.  This
includes both valuation-ring classes in
\texttt{LevelStructureConditions}.
\end{proposition}
\begin{proof}
The abelian locus is open, and its family is an abelian scheme.
For $(n,p)=1$ its $n$-torsion is finite \'etale.  A section on the
dense generic open selects a generic component with the same
function field as the normal base.  Its finite normal closure is
the base itself, giving the unique extension.  Applying this on
each component extends the original labels and every relation
between them.  This is precisely \texttt{Extends} for those labels.

For the second assertion, let $V$ be the specified valuation
ring and $K=\operatorname{Frac}V$.  Properness extends the generic
point uniquely to $y_V:\operatorname{Spec}V\to Y$.
Write $\mathcal G_V$ for the pulled-back semiabelian family and
$\mathcal B_V$ for the supplied abelian scheme.  The hypothesis
gives a group isomorphism between their generic fibres.  We
extend this particular isomorphism, without replacing $K$.

Choose an affine open of $Y$ containing the image of the closed
point of $\operatorname{Spec}V$.  It contains the entire image,
since every point of this local spectrum is a generization of
its closed point.  Its coordinate ring has image a finite-type
algebra over the original excellent arithmetic DVR $O$, inside
$V$.  Over any larger such subalgebra the pullback of the actual
family is still semiabelian.

Finite-presentation approximation descends $\mathcal B_V$, its
group operations and their identities to a finitely generated
$O$-subalgebra $R_1\subset V$.  Enlarge $R_1$ so that this
model is proper and smooth.  The locus of geometrically
connected fibres is open and contains the image of the closed
point of $V$; localize at one element outside that point, whose
inverse still belongs to $V$.  The descended group is then an
abelian scheme.  These are exactly the limit and connected-fibre
results of \cite{level-symmetry-approximation}.

The generic isomorphism and its inverse involve finitely many
coefficients in $K$.  Write them as fractions of elements of
$V$, and adjoin these numerators and denominators.  Finite
presentation, including the inverse and group identities, then
defines the two maps over $R_2[1/d]$ for a finite-type
$O$-domain $R_2\subset V$ and a nonzero $d\in R_2$.
Normalize $R_2$ in its fraction field.  The result $R$ is finite
over $R_2$, since $O$ is excellent.  It is contained in $V$:
every element of $R$ lies in $K$ and is integral over the
integrally closed valuation ring $V$.

We have thus obtained two semiabelian schemes over the normal
noetherian domain $R$, with the required inverse isomorphisms
on the dense open $D(d)$.  Proposition~3.3.1.5 of
\cite{level-symmetry-lan} extends both maps over $R$.
Uniqueness makes their composites identities.  Pulling back to
$V$ identifies $\mathcal G_V$ with $\mathcal B_V$, so every
point in the image of $y_V$ lies in the abelian locus.

A unit ball in the original nonarchimedean field is a valuation
ring with that fraction field.  Thus the argument proves the
nondiscrete clause as well as the complete-DVR clause, with all
original maps retained.  Completeness and discreteness are not
needed by this extension argument.
\end{proof}

These results, together with the whole-model integral theta
theorem, supply every field of \texttt{LevelStructureConditions}.
The simultaneous finite transitions and \texttt{SiegelCompare},
and the global domination clause, remain separate.
In particular a projective morphism between toroidal models is
not automatically the finite transition required by the source;
one must control its cone decomposition and integral normalization.
The remaining contracts stay explicit\cite{level-symmetry-source}.
